\subsection{A within-building design}
\label{sec:within-building}

Every estimate to this point identifies the text effect under an assumption, that
the location basis captures the spatial confounder, and bounds what an unobserved
violation could do. A subset of the data lets us replace that assumption with a
design. Units in the same building share location exactly. They sit at the same
coordinates, in the same neighborhood, the same school catchment, and the same
amenity set, so any confounder that is a function of location is constant across
them. Estimating the effect from variation in listing language between units of the
same building therefore differences spatial confounding out by construction rather
than adjusting it away, and what remains is identified by a design that holds
geography fixed.

The data support this on a scale worth using. Grouping listings by coordinates
rounded to about eleven meters, roughly eighteen thousand listings sit in buildings
that carry two or more, and the identifying variation is concentrated where it is
most useful, with New York contributing 6{,}258 listings across 1{,}885 buildings.
We add a building fixed effect, which is to say we demean the treatment, log price,
and the structural controls within building, and estimate the partially linear DML
coefficient on the within-building variation with standard errors clustered at the
building. Location controls are unnecessary here because the fixed effect absorbs
them.

The effect survives the design. Pooled across the twelve markets the within-building
coefficient is $+0.035$ ($t=10.8$), every market's point estimate is positive, and
the effect is significant in six of the twelve, among them New York at $t=8.6$,
Chicago at $t=8.6$, and Boston at $t=4.9$. The six that do not reach significance
are the markets with the fewest within-building listings, between roughly three
hundred and fifteen hundred, and each remains positive, so the pattern reads as
limited power rather than as a sign that geography was carrying the effect. New York
is the case that matters most. Its leading-component effect in the cross section was
an insignificant $+0.006$, and holding the building fixed returns a clean $+0.045$
at $t=8.6$, the same recovery the identified-direction estimator produced by a
different route. Two deconfounding strategies that share no machinery, one that
projects the location channel out of the text direction and one that differences the
building out of the design, agree that the New York effect is real and that the
variance-maximal direction had been hiding it.

Two limits keep this a strong check rather than a clean experiment. The building
fixed effect differences out everything that varies at the building level, which
removes the paper's primary threat but not a within-building one: a better-appointed
unit may draw both richer language and a higher price, and while the structural
controls for bedrooms, bathrooms, area, lot, and vintage absorb the coarse part of
that, they do not absorb unmeasured interior quality. The design identifies against
spatial confounding, not against unit quality. And the building key is coordinate
rounding rather than a parcel identifier, so some genuine buildings are split and
some neighbors are merged; both forms of grouping error attenuate the coefficient
toward zero, which makes the surviving effect a conservative reading. What the design
establishes is the claim the rest of the paper reaches by assumption, that a text
effect riding on a representation that encodes location is not reducible to that
location, now shown where location cannot vary at all.
